\documentclass[aps,pra,onecolumn]{revtex4-2}
\usepackage{amsmath,amssymb,booktabs}
\begin{document}

\title{Supplemental Material: Capacity-Normalized Bound on Information Backflow}

\maketitle

\section{Complete Proof with Domain Conditions}

\subsection{Definitions}

$N_S$, $N_E$: total qubits in system $S$ and environment $E$. The pointer basis $\{|0\rangle,|1\rangle\}$ is einselected by $H_{\rm int}$ [Eq.~(1), main text]. $q_S = N_S^{-1}\sum_{i\in S}\langle 0|\rho_i|0\rangle$, $q_E$ defined analogously. $C_F \equiv N_E q_E$ is the forward transfer capacity---the number of $E$-qubits initially available to receive a toggle. $N_{\rm back}$ is the number of backflow events ($E \to S$ toggles). $\mathcal{R} \equiv N_{\rm back}/C_F$ is the capacity-normalized backflow ratio.

\subsection{Proof}

\textbf{Step 1: Forward capacity.} The Lindblad jump $L_{S\to E}=|1\rangle\langle 1|_S\otimes|1\rangle\langle 0|_E$ toggles an $E$-qubit from $|0\rangle$ to $|1\rangle$ when the source $S$-qubit is $|1\rangle$. An $E$-qubit already in $|1\rangle$ is unaffected. The forward transfer capacity---the maximum number of $E$-qubits that can receive a toggle---is therefore $C_F = N_E q_E$.

\textbf{Step 2: Backflow bound.} Backflow events toggle $S$-qubits from $|0\rangle$ to $|1\rangle$ via $L_{E\to S}$. Since no reverse channel exists (neither $H_{\rm int}$ nor any Lindblad operator contains $|0\rangle\langle 1|$), each $S$-qubit can toggle at most once. The total number of $S$-qubits that can toggle during backflow is bounded by the number of $S$-qubits initially in $|0\rangle$: $N_{\rm back} \leq N_S q_S$ (in expectation).

\textbf{Step 3: Per-run vs.\ ensemble-average bound.} The integer number $n_0$ of $S$-qubits initially in $|0\rangle$ is a binomial random variable with mean $N_S q_S$. The strongest form of the bound is per-run and deterministic: $N_{\rm back} \leq n_0$, obtained by measuring $n_0$ at initialization via projective readout. No concentration inequality is needed for this formulation---it follows directly from the pigeonhole principle conditional on the measured $n_0$.

We caution against applying standard concentration inequalities (Hoeffding, Bernstein) to bound $N_{\rm back}$ directly. Toggle events are \textit{not independent}: each $|0\rangle \to |1\rangle$ toggle removes one qubit from the pool of available $|0\rangle$ states, introducing negative dependence between successive events. The sequence $\{X_t\}_{t=1}^k$ where $X_t = 1$ if a toggle occurs at step $t$ is a martingale difference sequence with respect to the filtration of the collision history, and the Azuma-Hoeffding inequality could in principle provide a concentration bound, but the per-run deterministic bound $N_{\rm back} \leq n_0$ is both simpler and strictly stronger. We therefore recommend the per-run formulation for experimental verification.

\textbf{Step 4: Ratio bound.} $\mathcal{R} \equiv N_{\rm back}/C_F \leq N_S q_S/(N_E q_E)$. For $N_S = N_E$: $\mathcal{R} \leq q_S/q_E$. $\square$

\subsection{Domain conditions}

The bound requires:
\begin{enumerate}
\item $q_E > 0$. If $q_E = 0$, all $E$-qubits are initially in $|1\rangle$, $C_F = 0$, and $\mathcal{R}$ is undefined.
\item $N_S, N_E \geq 1$ (non-empty system and environment).
\item The pointer basis is binary with unidirectional dynamics ($|0\rangle\langle 1|$ absent from $H_{\rm int}$ and all Lindblad operators).
\end{enumerate}

The bound is non-trivial ($\mathcal{R} < 1$) when $N_S q_S < N_E q_E$. For equal sizes, this is $q_S < q_E$, a regime accessible via state preparation.

\subsection{Why $\mathcal{R}$ is not the actual reflux fraction}

The actual fraction of forward-transfer events that reverse is $f_{\rm reflux} \equiv N_{\rm back}/N_F$, where $N_F \leq C_F$ is the actual number of forward events. Since $N_F \leq C_F$:
\begin{equation}
f_{\rm reflux} = \frac{N_{\rm back}}{N_F} \geq \frac{N_{\rm back}}{C_F} = \mathcal{R}.
\end{equation}
A bound on $\mathcal{R}$ provides a \textbf{lower} bound on $f_{\rm reflux}$, not an upper bound. To upper-bound $f_{\rm reflux}$, one would need a lower bound on $N_F$, which requires dynamical assumptions beyond the counting argument.

\section{Physical Origin of $L_{S\to E}$}

The jump operator $L_{S\to E}=|1\rangle\langle 1|_S\otimes|1\rangle\langle 0|_E$ describes a conditional excitation: when the source is $|1\rangle$ and the target $|0\rangle$, the target flips to $|1\rangle$. This operator arises from a microscopic model as follows.

Consider a resonant Jaynes-Cummings interaction between qubit $S$, qubit $E$, and a mediating bosonic mode:
\begin{equation}
H_{\rm micro} = g(\sigma_+^{(S)}\sigma_-^{(E)}a + \sigma_-^{(S)}\sigma_+^{(E)}a^\dagger),
\end{equation}
where $\sigma_\pm^{(i)}$ are Pauli ladder operators and $a$ is the mode annihilation operator. A dispersive measurement of $S$ in the pointer basis projects $\sigma_+^{(S)}\sigma_-^{(E)}$ onto its $|1\rangle\langle 1|_S$ component, yielding the conditional form. Tracing out the mode in the Born-Markov limit (weak coupling, short correlation time of the bath) gives the effective Lindblad operator $L_{S\to E}$.

The operator contains $|1\rangle\langle 0|$ but not $|0\rangle\langle 1|$. This asymmetry is energetic: the $|1\rangle$ state lies higher in energy than $|0\rangle$ by $\Delta E$. Spontaneous emission $|1\rangle \to |0\rangle$ occurs at a rate proportional to the electromagnetic density of modes at frequency $\Delta E/\hbar$, weighted by the Boltzmann factor $e^{-\Delta E/k_B T}$. When $\Delta E \gg k_B T$, this factor is exponentially small, and the reverse channel is suppressed. Typical superconducting qubit parameters ($\Delta E/h \sim 5$ GHz, $T \sim 15$ mK) give $e^{-\Delta E/k_B T} \sim 10^{-7}$.

\section{Time-Dependent $T_1$ Correction}

When $T_1$ relaxation is non-negligible, the population of $S$-qubits in $|0\rangle$ evolves during the protocol. Let $q_S(t)$ be the time-dependent fraction. The GKSL equation with a $|1\rangle \to |0\rangle$ decay channel at rate $\gamma = 1/T_1$ gives:
\begin{equation}
\frac{d}{dt} q_S(t) = \gamma(1 - q_S(t)).
\end{equation}
With initial condition $q_S(0) = q_S$, the solution is $q_S(t) = q_S + (1-q_S)(1-e^{-\gamma t})$. The number of $S$-qubits in $|0\rangle$ at protocol time $t$ is therefore:
\begin{equation}
N_S q_S(t) = N_S q_S + N_S(1-q_S)(1-e^{-\gamma t}).
\end{equation}
The second term represents $S$-qubits that were initially in $|1\rangle$ but decayed to $|0\rangle$, becoming available as backflow targets. The counting bound becomes $N_{\rm back} \leq N_S q_S(\tau_{\rm prot})$, where $\tau_{\rm prot}$ is the total protocol duration. The corrected bound is:
\begin{equation}
\mathcal{R} \leq \frac{q_S}{q_E} + \frac{(1-q_S)(1-e^{-\tau_{\rm prot}/T_1})}{q_E}.
\end{equation}
For $\tau_{\rm prot} \ll T_1$, expanding the exponential gives $\delta \sim (1-q_S)\tau_{\rm prot}/(q_E T_1)$, which is the form quoted in the main text.

\section{Relation to the MSV Entropic Bound and Budini Class}

Megier, Smirne, and Vacchini (MSV)~\cite{Megier:2021} proved that non-Markovian backflow is bounded by the Holevo skew divergence. The expression below is schematic only; the complete bound appears as Eq.~(4) of Ref.~\cite{Megier:2021}:
\begin{equation}
I_\mu(t,s) \leq \kappa\Big(\sqrt[4]{S_\mu(\varrho(s), \varrho_S(s) \otimes \varrho_E(s))} + \cdots\Big)
\end{equation}
(Schematic only; for the complete expression see Eq.~(4) of Ref.~\cite{Megier:2021}.)
This abbreviated form serves to indicate the qualitative structure---bound in terms of system-environment correlation---for comparison with the present work.
This bound applies to any CPTP dynamics and requires reconstruction of $\varrho_{SE}$, $\varrho_S$, and $\varrho_E$ via quantum state tomography. Our bound $\mathcal{R} \leq q_S/q_E$ constrains toggle-event counts rather than information-theoretic quantities, and the right-hand side requires only projective measurements.

Budini~\cite{Budini:2018} demonstrated that maximally non-Markovian quantum dynamics can exist without any environment-to-system backflow of information. In Budini's framework, non-Markovianity arises from time-dependent decoherence rates without physical information return from the environment. The present bound provides no constraint on this class: $N_{\rm back} = 0$ by construction, so $\mathcal{R} = 0 \leq q_S/q_E$ is trivially satisfied regardless of the degree of non-Markovianity. The bound is informative only when backflow is of the toggle type---a subclass of NM dynamics that includes collision models~\cite{Breuer:2016} and conditional excitation transfer in spin-boson systems.

\section{Gate-Level Considerations for Implementation}

\subsection{Stinespring implementation}

The Lindblad jump $L_{S\to E} = |1\rangle\langle 1|_S \otimes |1\rangle\langle 0|_E$ can be implemented via Stinespring dilation on $\{S, E, \text{ancilla}\}$: a Toffoli gate controlled by $\{S, \text{ancilla}\}$ targeting $E$, with the ancilla encoding whether $E$ is in $|0\rangle$. On superconducting hardware, the Toffoli decomposes into 6 CNOTs plus single-qubit gates~\cite{Shende:2009}. With CNOT fidelity $\sim 99.7\%$, a single Stinespring block has fidelity $\approx 98.2\%$.

For a protocol with $k$ Stinespring blocks, the cumulative gate fidelity is $(0.982)^k$. For $k=18$ (9 forward + 9 backflow on a $3 \times 3$ all-to-all system), this is $\approx 72\%$. The $\sim 28\%$ gate failure probability contributes to the systematic uncertainty on $\mathcal{R}$.

\subsection{Cross-resonance alternative}

A lower-depth approach uses the native $ZZ$ interaction on fixed-frequency transmons: a cross-resonance pulse $CR(\theta)$ on $\{S, E\}$ implements a controlled rotation conditioned on $S$ being in $|1\rangle$. This reduces the effective two-qubit gate depth from 6 CNOTs to $\sim 2$ cross-resonance pulses per block. Independent randomized benchmarking on the same qubits provides the gate fidelity without circular dependence on the bound measurement.

\subsection{Measuring $N_{\rm back}$}

The left-hand side of the bound requires counting $E \to S$ toggle events. On current superconducting hardware, this is achievable via mid-circuit measurements: after each collision step, the $S$-qubits are projectively measured in the pointer basis to determine whether a toggle occurred. The measurement collapses the state, freezing further dynamics---this is acceptable because each qubit toggles at most once, so a toggle event is irreversible and subsequent dynamics are independent of prior toggles. The total number of toggles across all $S$-qubits and all steps is $N_{\rm back}$.

\subsection{Distinguishing gate error from bound violation}

If a measured $\mathcal{R}$ exceeds $q_S/q_E$, three explanations are possible: (a) gate error (a qubit toggled twice), (b) $T_1$ decay inflating the effective $q_S$ beyond the corrected estimate, or (c) breakdown of the unidirectional-toggle picture. Independent gate characterization---interleaved randomized benchmarking before and after the protocol---bounds the gate-error contribution. If the excess $\mathcal{R}$ exceeds the $3\sigma$ upper limit from gate characterization and the $T_1$ correction $\delta$, the unidirectional-toggle model is excluded.

\begin{thebibliography}{3}

\bibitem{Breuer:2016}
H.-P.~Breuer, E.-M.~Laine, J.~Piilo, and B.~Vacchini,
Colloquium: Non-Markovian dynamics in open quantum systems,
Rev.\ Mod.\ Phys.\ \textbf{88}, 021002 (2016).

\bibitem{Megier:2021}
N.~Megier, A.~Smirne, and B.~Vacchini,
Entropic bounds on information backflow,
Phys.\ Rev.\ Lett.\ \textbf{127}, 030401 (2021).

\bibitem{Budini:2018}
A.~A.~Budini,
Maximally non-Markovian quantum dynamics without environment-to-system backflow of information,
Phys.\ Rev.\ A \textbf{97}, 052133 (2018).

\bibitem{Shende:2009}
V.~V.~Shende and I.~L.~Markov,
On the CNOT-cost of Toffoli gates,
Quantum Inf.\ Comput.\ \textbf{9}, 461 (2009).

\end{thebibliography}

\end{document}
